\section{Proprioceptive embodiment adaptation}
\label{sec:vla-method}
The evaluated policies are from the $\pi_0/\pi_{0.5}$, OpenVLA-OFT, and
GR00T families \citep{pi0,pi05,openvlaoft,grootn1,grootn17}.
\subsection{Execution interface and calibration}
At physical command step $k$, let $a_k$ be the nominal action decoded by the
retained VLA and $D$ the fixed diagonal correction mask. Let
$\fhat_k^{\rm app}=\fhat_{\tau_k}$ be the estimate snapshot used for this
command, with $0\le\tau_k\le k$. The execution map includes the robot
and its low-level controller. At the modeled interface, write
\begin{equation}
u_k=a_k-D\fhat_k^{\rm app}+\zeta_k,\qquad
u_k^{\rm eff}=u_k+f_k.
\label{eq:vla-command}
\end{equation}
Here $f_k$ is an additive physical offset in the interface's command units.
The discrepancy $\zeta_k$ includes imperfect realization of the intended
correction and any clipping relative to that correction. Exact external
subtraction with no clipping has $\zeta_k=0$; editing a native decoder bias
generally does not: a registered diagnostic on the native
\code{action\_out\_proj/bias} edit measures a median realization error of
43\% of the intended correction, because the decoder gain varies 27--41\%
across observations (Appendix~\ref{app:gen-authority}). Every published
cohort subtracts the correction outside the network.

Fit a healthy finite-impulse-response (FIR) predictor, a sensitivity matrix
$S$, and a healthy reference $b_{\rm cal}$ in inverse-sensitivity coordinates:
\begin{equation}
\widehat y_k=\sum_{\ell=0}^{K}H_\ell u_{k-\ell}+\widehat c,
\qquad e_k^{\rm res}=y_k-\widehat y_k,\qquad
z_k=S^{-1}e_k^{\rm res}-b_{\rm cal}=f_k+w_k.
\label{eq:vla-observation}
\end{equation}
The response $y_k$ is acquired after executing $u_k$ and updates the next
estimate. The predictor uses the command actually presented at the modeled
interface, before the unknown fault, with response timestamps aligned to that
command. An internal intended correction is not interchangeable with the
final decoder/clipping output. Healthy rollouts fit $H_\ell,\widehat c$;
controlled offset probes identify $S$. The error $w_k$ contains history,
model/sensitivity mismatch, noise, and imperfect healthy centering
(Appendix~\ref{app:centering}). The issued-command FIR buffer is observer
memory; physical/controller delays remain in the execution state, while VLA
internal memory belongs to the command source (Appendix~\ref{app:history-partition}).



\begin{table}[H]
\centering\small
\caption{Documented interface families. New robots receive new calibration.
The listed update families follow the documented mathematical descriptions.}
\label{tab:implementation}
\setlength{\tabcolsep}{4pt}
\begin{tabular}{@{}>{\raggedright\arraybackslash}p{.16\linewidth}>{\raggedright\arraybackslash}p{.30\linewidth}>{\raggedright\arraybackslash}p{.23\linewidth}>{\raggedright\arraybackslash}p{.22\linewidth}@{}}
\toprule
Interface & Predictor response & Offset estimate & Evaluation schedule\\
\midrule
LIBERO Panda & Position increments; relative rotations ($K=6$) & Observation attenuation & Within episode\\
ALOHA & Absolute joint positions ($K=6$) & Observation attenuation & Prior identify, then hold\\
Fourier GR1 & Absolute joint positions ($K=6$) & Innovation normalization & Three prior episodes; held median\\
\bottomrule
\end{tabular}
\end{table}

\subsection{Documented update and causal timing}
Let $P_r$ select residual coordinates in the run's declared units. Its
normalization scale $\rho$ is tied to those coordinates; changing units or
whitening changes the estimator. Position increments and relative-rotation
coordinates have different physical units; exact scaling must accompany each
configuration. $P_r$ is distinct from
$D$: an uncorrected coordinate can still alter the shared gain; the
recovered configurations use all six residual coordinates on LIBERO and all
fourteen on ALOHA, while GR1 omits its unmodeled hand coordinates. LIBERO and ALOHA use observation attenuation,
while GR1 uses innovation normalization:
\begin{align}
\mathcal T_k^{\rm att}(\fhat_k)
&=\Pi_{\mathcal C}[(1-\gamma)\fhat_k+\gamma s_kz_k],
&s_k&=\frac{\mathbf1[\|P_re_k^{\rm res}\|>\delta]}{1+\|P_re_k^{\rm res}\|^2/\rho^2},
\label{eq:legacy}\\
\mathcal T_k^{\rm inn}(\fhat_k)
&=\Pi_{\mathcal C}[\fhat_k+\alpha_k(z_k-\fhat_k)],
&\alpha_k&=\frac{\gamma}{1+\|P_rS(z_k-\fhat_k)\|^2/\rho^2}.
\label{eq:innovation}
\end{align}
The set $\mathcal C$ bounds the estimate; the selected-coordinate results use a coordinatewise box. For a full-update flag
$\chi_k\in\{0,1\}$, apply
\begin{equation}
\fhat_{k+1}=\begin{cases}\mathcal T_k(\fhat_k),&\chi_k=1,\\
\fhat_k,&\chi_k=0.\end{cases}
\label{eq:method-schedule}
\end{equation}
A zero observation gate in Eq.~\eqref{eq:legacy} retains leakage when
$\chi_k=1$; $\chi_k=0$ holds the entire estimate. The analysis treats both
schedules explicitly. The snapshot $\tau_k$ separately captures delayed
application to an action chunk, even if estimation is updated more often.
LIBERO reference settings are $\gamma=.08$, $\rho=.15$, $\delta=.008$,
and a $\pm.15$ estimate box in action coordinates.

\begin{algorithm}[H]
\centering\small
\begin{minipage}{.97\linewidth}
\textbf{Inputs:} calibrated predictor and sensitivity; residual coordinates;
mask $D$; projection set; update family; gate and snapshot-application rule.
\begin{enumerate}[leftmargin=1.5em,itemsep=1pt,topsep=3pt]
\item Obtain the next nominal action from the decoder or buffered chunk;
select the available snapshot $\fhat_{\tau_k}$.
\item Realize the correction, apply limits, and execute. Save the actual
command $u_k$ at the modeled interface and its timestamp.
\item Associate the next proprioceptive response with that command and form
$e_k^{\rm res},z_k$ from the synchronized history and healthy reference.
\item Set $\chi_k$ using the declared schedule, update by
Eq.~\eqref{eq:method-schedule}, and store the estimate for later application.
\end{enumerate}
\end{minipage}
\caption{Causal execution semantics used in the analysis. Fault truth is
unavailable to the deployment estimator. The observed update/hold and
application schedules must be logged to instantiate the schedule bounds.}
\label{alg:execution}
\end{algorithm}

\subsection{Correction support, gain faults, and evidence scope}
The primary LIBERO condition estimates six coordinates and applies three
rotation corrections. Translation/rotation masks were informed by fault-family
diagnostics; the evaluations therefore use family-informed support and do
not establish autonomous unknown-support selection. Three healthy episodes
(285 steps) calibrate LIBERO; that calibration is reused across policies.
ALOHA and GR1 use eight and six healthy episodes, respectively. Probe and
fault-development data are additional to those healthy budgets.

For a matched diagonal FIR with constant actuator gain, define
$\phi_{i,k}=\widehat y_{i,k}-\widehat c_i$ and
$e_{i,k}^{\rm res}=(g_i-1)\phi_{i,k}+\varepsilon_{i,k}$. The gain branch uses gated
recursive least squares and feasible inverse correction $a_i/\widehat g_i$.
Gains must remain bounded away from zero; joint gain/offset estimation
additionally needs informative regressor variation.

A registered provenance pass closed the remaining implementation
questions: every published cohort subtracts the correction outside the
network, so native-edit realization error never enters the reported cells;
the gate is a leakage deadzone on LIBERO and ALOHA and a full hold on GR1,
and it never fired in the four headline samples or ALOHA's identification
episodes; $K=6$ throughout. The realization and schedule terms remain in the
analysis because the gated and native-edit regimes the map documents need
them (Appendix~\ref{app:experiment}).
